% =====================================================================
%  The Shape of the Threat — front matter
%  Half-title, title page, copyright, dedication, epigraph, contents,
%  figures, tables, management summary, preface, how to use this book.
% =====================================================================

% ---- Half-title ---------------------------------------------------------
\pagestyle{plain}
\clearpage
\thispagestyle{empty}
\vspace*{1.15in}
{\sffamily\fontsize{22}{26}\selectfont\color{txt}\bookTitle\par}
\vspace{10pt}
{\sffamily\fontsize{11}{14}\selectfont\color{txtdim}\bookSubtitle\par}
\vfill
{\sffamily\small\color{txtdim}\bookAuthor\par}
\vspace{6pt}
{\sffamily\small\color{txtdim}\bookImprint\par}
\clearpage

% ---- Title page ---------------------------------------------------------
\clearpage
\thispagestyle{empty}
\vspace*{0.75in}
{\sffamily\fontsize{26}{30}\selectfont\bfseries\color{txt}\bookTitle\par}
\vspace{12pt}
{\sffamily\fontsize{12}{16}\selectfont\color{accent}\bookSubtitle\par}
\vspace{22pt}
\noindent{\color{accent}\rule{\textwidth}{1.8pt}}\par
\vspace{18pt}
{\sffamily\large\color{txt}\bookAuthor\par}
\vspace{8pt}
{\sffamily\small\color{txtdim}\bookImprint, publisher\par}
\vspace{4pt}
{\sffamily\small\color{txtdim}\bookEdition\par}
\vspace{18pt}
{\sffamily\small\color{txtdim}
Facts current as of \bookFactsDate.\par}
\vfill
\clearpage

% ---- Copyright page -----------------------------------------------------
\clearpage
\thispagestyle{plain}
\noindent{\sffamily\bfseries\color{txt}\bookTitle}\par
\noindent{\sffamily\small\color{txtdim}\bookSubtitle}\par
\vspace{14pt}
\noindent\rule{\textwidth}{0.4pt}\par
\vspace{10pt}
\noindent Copyright \textcopyright\ \bookCopyright\ \bookAuthor. All rights reserved.\par
\vspace{6pt}
\noindent \bookImprint, publisher.\par
\vspace{10pt}
\noindent\textbf{Edition:} \bookEdition.\par
\vspace{6pt}
\noindent\textbf{Subject headings:} Near-Earth asteroids --- Observational
techniques; Planetary defense; Asteroid shape; Impact energy; Torino scale;
Interferometry; Optical ranging; Lunar surface instrumentation; Space mission
planning.\par
\vspace{10pt}
\noindent\textbf{Licence.} This work is published under the Apache License,
Version 2.0. The complete licence text is available at
\href{https://www.apache.org/licenses/LICENSE-2.0}{\code{apache.org/licenses/LICENSE-2.0}}.
The book may be copied, modified and redistributed, including for commercial
purposes, provided that this notice and the attribution required by the
licence are retained.\par
\vspace{10pt}
\noindent\textbf{DOI:} \href{https://doi.org/\bookDOI}{\texttt{\bookDOI}}\par
\vspace{10pt}
\noindent\textbf{ISBN:} 978-1-0000-0000-0 (placeholder; to be assigned on
publication).\par
\vspace{10pt}
\noindent\textbf{Source materials.} The sources of this book, the build
scripts, the figures and the research ledger behind it are published at
\texttt{github.com/lordxmen2k/The-Shape-of-the-Threat}.\par
\vspace{10pt}
\noindent\textbf{What this book is.} This book is an argument about where a
near-Earth object characterisation capability should sit inside the
planetary-defence architecture, and about what it would cost and buy. It works
from primary sources, carries its own quantitative analysis, and states the
limits of that analysis plainly. It does not describe a spacecraft that has
been flown, a station that has been built, or a decision that has been taken.
It does not cover deflection, and it does not cover the economics of a
planetary-defence programme.\par
\vspace{10pt}
\noindent\rule{\textwidth}{0.4pt}\par
\vspace{8pt}
{\sffamily\small\color{txtdim}Typeset in XeLaTeX with Outfit, DejaVu Sans Mono
and Fira Math. First edition, 2026.\par}
\clearpage

% ---- Dedication ---------------------------------------------------------
\clearpage
\thispagestyle{empty}
\vspace*{2.6in}
\begin{center}
{\sffamily\itshape\color{txtdim}
For the objects that were found in time to be looked at,\\[6pt]
and were not looked at.\par}
\end{center}
\vfill
\clearpage

% ---- Epigraph -----------------------------------------------------------
\clearpage
\thispagestyle{plain}
\begin{epigraph}
No automated system currently alerts physical characterization sites.

\rule{1.6in}{0.3pt}

NASA, \emph{Lessons Learned from the 2024 YR4 Impact Scenario}, 2025
\end{epigraph}
\vfill
\clearpage

% ---- Contents -----------------------------------------------------------
\clearpage
\thispagestyle{plain}
\markboth{Contents}{Contents}
\tableofcontents
\addcontentsline{toc}{chapter}{Contents}
\clearpage

% ---- Figures ------------------------------------------------------------
\clearpage
\thispagestyle{plain}
\markboth{Figures}{Figures}
\listoffigures
\addcontentsline{toc}{chapter}{List of Figures}
\clearpage

% ---- Tables -------------------------------------------------------------
\clearpage
\thispagestyle{plain}
\markboth{Tables}{Tables}
\listoftables
\addcontentsline{toc}{chapter}{List of Tables}
\clearpage

% ---- Management summary -------------------------------------------------
% The front matter's unnumbered chapters must not number their sections.
\setcounter{secnumdepth}{-1}
\chapter*{Management Summary}
\addcontentsline{toc}{chapter}{Management Summary}
\markboth{Management Summary}{Management Summary}
\thispagestyle{plain}

This page is the book in the form a decision-maker needs it. The remainder of
the volume derives, tests and qualifies every statement below.

\section*{The question}

The near-Earth object catalogue contains more than forty thousand objects. A
minority of them carry a physical account good enough to size them. The
planetary-defence community can therefore say, for most known objects, roughly
how much energy would be released on impact --- and only roughly. The
question this book asks is where in the planetary-defence architecture a
measurement capability should be placed so that the objects that matter are
accounted for rather than merely catalogued.

\section*{The answer}

A fixed surface station on the Moon, carrying a kilometre-scale optical
interferometric array, operated as a \emph{named-target characterisation
service} rather than as a survey. It is not the only answer, and on two of the
five jobs it loses to platforms that already exist. It is the answer for the
one job that nothing else is assigned to do: measuring the resolved shape of a
specific object, at the moment when that object's size decides whether an
impact is regional or local.

\section*{The rule}

For a target with a bracketed diameter, density, distance and encounter epoch:
compute the baseline required for four samples across the diameter, which is
four wavelengths divided by the angular diameter; check the baseline is
buildable at the candidate site; check the astrometric precision implied by
that baseline against the along-track budget at the encounter horizon; and
select the platform only if the \emph{physical} uncertainty dominates the
encounter probability. That last condition is the substantive one. Where the
astrometric term dominates, characterisation does not shorten the clock, and
the lunar station does not earn its place.

\section*{The three numbers that decide it}

\begin{itemize}
\item \textbf{748 metres.} The baseline a 60~metre object at 0.1 astronomical
units requires for four samples across its diameter. The concept literature's
100~metre baseline puts a 50~metre object at the same distance at 0.45 of a
fringe spacing across its diameter: an astrometric centroid, not a shape
imager.
\item \textbf{8.66 metres.} The range error of an Earth-Moon two-station
measurement of 2024~YR4 at the distance of its discovery-night close approach,
at one milliarcsecond of differential astrometry. No ground optical facility
ranges a non-cooperative target.
\item \textbf{37 campaigns per year.} The service rate of one fixed station,
against 342 new 140-metre-class objects a year and 11,725 known. The station
cannot be a census. It must be a triage.
\end{itemize}

\section*{What it does not do}

It does not find objects. It does not carry the thermal-infrared channel,
which Sun-Earth L1 does better and more cheaply. It does not survey. It does
not deflect. It does not clear the catalogue: at 37 campaigns a year, covering
the known 140-metre class would take 313 years, and the arithmetic says so
before any hardware is designed.

\begin{keytakeaways}
\item The gap for the hazard-relevant class is accounting, not discovery, and
the book shows the arithmetic that says so.
\item The lunar claim is narrowed to the only thing the Moon does that nothing
else is assigned to do, and the claim is made quantitative rather than
rhetorical.
\item The platform-selection rule returns a platform per target, not an
architecture, and it declines to select the Moon whenever the astrometric term
dominates.
\item The service rate is stated against the demand, so the station is
proposed as a triage capability and not as a census.
\end{keytakeaways}
\clearpage

% ---- Preface ------------------------------------------------------------
\chapter*{Preface}
\addcontentsline{toc}{chapter}{Preface}
\markboth{Preface}{Preface}
\thispagestyle{plain}

A near-Earth object is a line in a catalogue: a designation, an orbit, an
absolute magnitude, and a probability of impact that is usually negligible. The
line carries no size, no mass, no shape and no spin state, because none of
those have been measured for most of the objects on it. That is not a criticism
of the surveys that built the catalogue. It is a statement of what a catalogue
is for.

The planetary-defence community has inherited a second problem from this one,
and it is the problem this book is about. The hazard scale that the community
uses to decide whether an object deserves attention has two axes: the
probability that the object strikes, and the energy it would release. The
first axis is measured, by an astrometric machinery that has grown
extraordinarily well. The second axis is inferred, from a size assumption and
a density assumption, and it is good to within a factor of three. A factor of
three in energy is the difference between an event that empties a regional
hospital and one that levels a city.

Shape sits at the head of both chains. The accounting chain runs shape to
volume to mass to energy to hazard class. The probability chain runs shape to
surface normals to YORP torque to spin state to Yarkovsky drift to along-track
displacement to keyhole. The two chains are usually presented separately, and
they are usually presented as separate programmes. They are the same
measurement, and the measurement is the same observation: a resolved image of
a rotating body.

\section*{What the book does}

It takes the accountability chain apart and attaches an error budget to every
link, so that the reader can see which measurement would buy the most. It
takes the probability chain apart and shows why the physical term is
negligible at seven years and dominant at thirty. It re-derives, from the
concept study's own published table, the baseline that an asteroid shape
imager actually requires, and finds that the baseline in the literature is
sized for stellar astrometry rather than for asteroid imaging. It then works
through the lunar case honestly, including the two channels where the Moon
loses, and closes with a rule that returns a platform for a given target
rather than a preferred architecture.

It also does something the literature does not do: it tests the proposal
against the historical record. The counterfactual chapter asks what a lunar
capability would have measured in the cases that are already on the record ---
2008~TC3, 2019~OK, 2024~YR4, Apophis --- and answers with the JPL close
approach and first-observation records rather than with an assertion.

\section*{What the book will not do}

It will not recommend a mission. The platform-selection rule is the
deliverable, and the rule is indifferent to the answer.

It will not claim novelty of hardware. The concept is twenty-three years old
and heavily worked. The novelty claimed here is of position and of
quantification: the first published argument for where in the
planetary-defence architecture a lunar platform earns its place, together with
the error budgets and the rule that decide it.

It will not treat a model as a measurement. Where the book uses a thermal
model, a beaming parameter or a scaling law, the error band is stated and the
places where the extrapolation does not follow are named.

\section*{On the sources}

Every identifier in this book was resolved against a registry --- the arXiv
API, the Crossref bibliographic query, the NASA Technical Reports Server
citation service, or the NASA/JPL Small-Body Database --- and the resolved
title was compared against the title intended to be cited before the reference
was used. The discovery statistics are the NASA Center for Near Earth Object
Studies series as published for 22 November 2026, and the close-approach
records are the JPL Small-Body Database Close Approach Data service. Where a
resolution failed, the claim was removed rather than re-sourced from memory.

The sources are institutional, peer-reviewed, or official. Where the evidence
is thin, the book says so, and where two sources disagree, both are carried
with their definitions rather than reconciled by preference.

\begin{keytakeaways}
\item Shape is the shared upstream term of the accounting chain and the
probability chain, and it is the book's centre of gravity.
\item The baseline an asteroid shape imager needs is a kilometre-scale
baseline, not a hundred-metre one, and the literature's hundred-metre figures
are sized for stellar astrometry.
\item The lunar claim survives only when it is narrowed to what the Moon does
that no space observatory is currently assigned to do.
\item Every identifier was resolved against a registry and checked against the
intended title before use.
\end{keytakeaways}

% ---- How to use this book ------------------------------------------------
\chapter*{How to Use This Book}
\addcontentsline{toc}{chapter}{How to Use This Book}
\markboth{How to Use This Book}{How to Use This Book}
\thispagestyle{plain}

The book has a spine, and reading it out of order is possible but wasteful.
The seven parts are a sequence, not a taxonomy.

\textbf{Part I} establishes the gap. What a catalogue entry cannot tell you,
the two axes of the hazard scale, the difference between the detection mandate
and the characterisation mandate, the size of the population that is discovered
but not accounted for, and 2024~YR4 as the case where the response failed.

\textbf{Part II} is the physics of accountability: what has to be measured to
turn a catalogue entry into an energy, and how much error each link carries.

\textbf{Part III} is the probability chain, where the same shape measurement
turns into a Yarkovsky drift and a keyhole probability, and where the
sensitivity of the answer to physical error grows with the horizon.

\textbf{Part IV} is the measurement science: projections and the uniqueness
theorems, lightcurve inversion and its convex hull, radar and the fourth-power
wall, resolved imaging and the close-approach window, and interferometry with
its visibility and over-resolution limits.

\textbf{Part V} is the Moon: what the vantage gives, what the baseline gives
and what it does not, the thermal and power budget, dust and night, and the
architecture and service model that follow.

\textbf{Part VI} is the comparison and the rule: the space telescopes, bistatic
radar, the portfolio decision, the counterfactual ledger against the
historical record, the platform-selection rule, the worked cases, and the
honest novelty claim.

\textbf{Part VII} states what a lunar station could and could not be made
responsible for, and what would falsify the argument.

A reader who wants the argument rather than the detail should read Chapters 1,
4, 10, 14, 20, 22, 28, 29 and 30, and skip the rest. A reader who is
commissioning a capability should read them in order and use the glossary and
the index.

\begin{table}[!htbp]
\centering
\small
\caption{The book's own computed figures, and what each one is for.}
\label{tab:figures}
\begin{tabularx}{\textwidth}{@{}P{2.57in}Y@{}}
\toprule
\rowcolor{panel}
\textbf{Figure} & \textbf{The question it answers} \\
\midrule
Known near-Earth objects & how large the catalogue is, and where its growth
has gone \\
\rowcolor{panel}
The 140 m class & how many objects are catalogued but not accounted for \\
Baseline for four samples & what baseline an asteroid shape imager actually
requires \\
\rowcolor{panel}
Samples across the diameter & where the concept literature's hundred-metre
baseline fails \\
The mid-infrared channel & why the thermal channel is closed to the Moon \\
\rowcolor{panel}
The accountability ladder & what each link of the energy chain is worth \\
Two-station range error & what the Earth-Moon baseline buys at close approach \\
\rowcolor{panel}
2024 YR4 & how far model choice alone moves the energy \\
The 2036 keyhole & how far a physical assumption moves the probability \\
\rowcolor{panel}
Along-track error & when the physical term overtakes the astrometric one \\
The sky turns more slowly & what the lunar rotation rate buys \\
\rowcolor{panel}
The Sun-ward lead time & what the lunar survey geometry adds \\
Baseline scales & what has been achieved against what is required \\
\rowcolor{panel}
Surface brightness & why the lunar sky is not the constraint on the optical
channel \\
\bottomrule
\end{tabularx}
\end{table}
